\chapter{Simulation 03: switching and recorded history}\label{ch:s03}
\plain{Each center now carries a switch as well as its position and velocity.
The switch reverses the new interaction contribution, not the motion already
being carried. The saved logs reveal the difference. They also show that
the program can restart its labelled clock without restarting its state.
This chapter follows those records before interpreting the movie.}

\section{Construction and added state}
The 167-line \filepath{plot_1.py} retains eight 3D centers, period~10,
$\Delta t=0.1$, the same tetrahedral offsets, random initial coordinates
and random velocity components. It adds eight binary modes $b_i$, all
initialized to zero. Its numerical carried state therefore has 48 real
entries and eight mode bits, together with the random stream and externally
supplied frame input. \src{S03}{10--28,152}; \cite{study03}.
\claim{C45}

The common H, componentwise RPCO and mean-speed GRF function bodies match
S02 after removing docstrings, as recorded by the accepted AST check.
At each callback, GRF and H are computed once, then the incoming positions
are copied for a diagnostic. Bodies update in order 0--7, with the same
mixed-state arithmetic centroid and wrapped displacement.
\src{S03}{36--54,102--113}. \claim{C46}

The source is not merely S02 plus a color change. It also changes nominal
frame count to 200, requests \filepath{blit=False}, saves at 20 fps,
computes a kinetic diagnostic, prints a jump diagnostic and appends
numerical records on every callback. These changes affect the available
chronology and presentation even though they do not all feed into the
movement rule. \src{S03}{12,57--78,129--167}. \claim{C47}

\section{Exact switch and update equations}
Use $\W_{10}$, $\D_{10}$ and $\R$ from
Equations~\eqref{eq:wrap10} and~\eqref{eq:RPCO}. At callback $q$,
with incoming velocities $v_j$ and supplied label $f_q$, set
\begin{equation}\label{eq:shared03}
g_q=\frac{1}{1+\frac18\sum_j\|v_j\|_2},\qquad h_q=H(0.1f_q).
\end{equation}
For each body $i$, calculate $c_{q,i},d_{q,i},r_{q,i}$ using
Equations~\eqref{eq:centroid}--\eqref{eq:centroiddiff}. Only then
evaluate the gate. If $g_q>0.5$, draw $U_{q,i}$ using
\filepath{np.random.rand()}; otherwise no draw is consumed.
Define
\begin{equation}\label{eq:toggle}
F_{q,i}=
\begin{cases}
1,&g_q>0.5\ \hbox{and the consumed draw }U_{q,i}<0.01,\\
0,&\hbox{otherwise},
\end{cases}
\qquad
b_i^+=b_i\mathbin{\mathrm{XOR}}F_{q,i}.
\end{equation}
The update is
\begin{align}
a_{q,i}
 &=0.1g_qh_q(1-2b_i^+)\frac{d_{q,i}}{r_{q,i}+0.01},\label{eq:kick03}\\
u_{q,i}&=v_i+a_{q,i},\qquad v_i^+=\R(u_{q,i}),\label{eq:vel03}\\
x_i^+&=\W_{10}(x_i+0.1v_i^+).\label{eq:pos03}
\end{align}
The sign is the \emph{post-toggle} mode, the one recorded in the current
block. H is multiplied once, and GRF remains the incoming common value
throughout the sweep. \src{S03}{102--126}. \claim{C48}

Conditional on the intended uniform generator and an open gate, the
toggle probability is 1\% per body per callback. No random values or RNG
state are logged. When H is positive, mode~0 contributes toward the
chosen wrapped centroid displacement and mode~1 away from it. When H is
negative those roles reverse. Neither mode alone specifies the outgoing
velocity direction. \claim{C49}

\section{Recording and callback chronology}
At startup the source writes an eleven-column header to a timestamped
CSV. After each complete callback it appends eight rows in body order:
\begin{center}\small
\filepath{Step,Tetrahedron,X,Y,Z,VX,VY,VZ,Energy,GRF,Rule_State}.
\end{center}
Positions, velocities and modes are outgoing values; Energy is the
whole-system post-sweep value repeated in each row; GRF is incoming.
There is no saved detector flag, H column, within-sweep snapshot or
per-row wall-clock time. \src{S03}{74--78,129,144--147}. \claim{C50}

\begin{table}[htbp]\centering\small
\caption{Original log structure, preserved in physical row order.
The repeated-key count is the number of rows beyond the first occurrence
of each (Step,Tetrahedron) key. Accepted source: S03
ORIGINAL\_LOG\_ANALYSIS and analysis/log\_summary.json.}\label{tab:counts03}
\begin{tabular}{@{}lrr@{}}\toprule
Measure & Log A & Log B\\\midrule
Data rows &8,736&3,520\\
Eight-body callback blocks &1,092&440\\
Distinct frame labels &200&200\\
Label-zero blocks &12&6\\
Repeated-key rows &7,136&1,920\\
Partial or corrupt blocks &0&0\\\bottomrule
\end{tabular}
\end{table}

Log A is \filepath{PLOT_1_tetrahedral_states_20250712_015421.csv};
log B is \filepath{tetrahedral_states_20250712_015647.csv}.
They contain finite values, correctly ordered body IDs and shared
Step/Energy/GRF within each block. All positions lie in $[0,10)$ and
all modes are binary. They are strongly attributed to this recurrence
by schema and numerical agreement; the original seed and pre-log
state are missing. \claim{C51}

\begin{table}[htbp]\centering\small
\caption{Exact label chronology. A dash denotes no further record,
not an assumed reset or simulated continuation.}\label{tab:segments03}
\begin{tabular}{@{}lll@{}}\toprule
Callback range & A labels & B labels\\\midrule
0 &0&0\\
1--200 &0--199&0--199\\
201 &0&0\\
202--401 &0--199&0--199\\
402 &0&0\\
403--602 &0--199&0--36 through q439 only\\
603 &0&---\\
604--803 &0--199&---\\
804 &0&---\\
805--1004 &0--199&---\\
1005 &0&---\\
1006--1091 &0--85&---\\\bottomrule
\end{tabular}
\end{table}

Each label-zero block is a distinct recorded state. Sorting by Step and
discarding duplicates would destroy the chronology in
Table~\ref{tab:segments03}. Counting $0.1$ per recorded callback gives
109.2 and 44.0 model-step units; the first-to-last stored-state spans
are 109.1 and 43.9. Neither count is the variable supplied to H:
its argument repeatedly returns to zero.
\cite{study03}, ORIGINAL\_LOG\_ANALYSIS. \claim{C52}

The accepted inspection of local Matplotlib 3.10.9 identifies
initialization, save, interactive start and repeated playback as
possible sources of extra first-label calls. The observed singleton
zeros and complete passes are consistent with those paths. However,
the historic library, backend and GUI event sequence are missing.
We preserve recorded labels without assigning an invented save or
window event to every segment. \claim{C53}
\EvidenceFig{F10}

\section{What the accepted numerical checks cover}
The accepted primary checker read every row. It reconciled the kinetic
quantity on all 1,532 blocks, and position and incoming-GRF relations
on every consecutive block: 8,728 body transitions in A and 3,512 in B.
Callback zero lacks a recorded predecessor, so its incoming update
cannot be certified. \claim{C54}

The full velocity equation was checked at a declared union of sample
callbacks: the first five transitions, final five, maximum native
increment, strongest consecutive-vector turn, closest pair, first
label discontinuity, first decreasing-label reset, first switch and
first seam. Deduplicating that \emph{selection list}, not the data,
gave 16 sampled transitions in A and 15 in B. The checker substituted
the observed current prefix and preceding suffix into the centroid;
it did not propagate calculated states. An independent scalar
implementation checked the 20 original R0 sample transitions and
all energy/position/GRF identities.
\cite{study03}, equation-residual CSVs and independent\_checks.json.
\claim{C55}

All declared checks passed an absolute tolerance of $10^{-12}$.
Across the methods, reported maximum position residual was
$1.78\times10^{-15}$ and sampled velocity residual
$3.61\times10^{-16}$. Such agreement supports source attribution and
the interpretation of columns. It is not an exhaustive velocity
audit, recovery of random draws or bit-identical reproduction of
the unknown historical execution. These scientific checkers were
not rerun for Volume I. \claim{C56}

\section{A sign change before movement reverses}
All recorded callbacks have GRF above 0.5. Its range is approximately
0.633309--0.987992 in A and 0.641949--0.987362 in B.
Between successive recorded blocks there are 79 toggles in A and
33 in B. All are gate-eligible, but the unrecorded draws cannot be
reconstructed or used to certify a distribution. \claim{C57}

The shorter log's first toggle is body~5 at $q=7$, Step~6, CSV line~63.
It changes mode from 0 to 1 while H remains positive.
Table~\ref{tab:delay03} follows the recorded response. The radial
component is the dot product of outgoing velocity with the unit
wrapped centroid displacement evaluated at that body's update,
using the observed mixed-state positions. Positive means along
that instantaneous direction. \claim{C58}

\begin{table}[htbp]\centering\small
\caption{Original log B, body 5. A response-sign change at q7 precedes
an outward radial velocity at q9. Source: original CSV lines shown;
accepted analysis/episode\_records.json. Rounded values, no new update.}
\label{tab:delay03}
\begin{tabular}{@{}rrrrr@{}}\toprule
Callback & CSV line & Mode & Speed & Radial outgoing velocity\\\midrule
6&55&0&0.570163&$+0.565607$\\
7&63&1&0.291844&$+0.284300$\\
8&71&1&0.087759&$+0.062268$\\
9&79&1&0.141476&$-0.131445$\\\bottomrule
\end{tabular}
\end{table}

The new opposing contribution first reduces an inward carried
velocity. The next record still moves inward; two callbacks after
the toggle, its radial outgoing component is negative. This is a
specific delayed direction change, not a universal two-callback
law. A's first toggle, body~4 at $q=12$, has an analogous ordering
through $q=14$. \claim{C59}

Across all observed toggles, none has a negative dot product
between the immediately preceding and outgoing velocity vectors.
There are 17 such obtuse consecutive turns in A and 6 in B at
other transitions. The strongest examples are A body~4 at
$q=1005$ (line 8046, cosine about $-0.977692$) and B body~5 at
$q=402$ (line 3223, about $-0.973052$). Both retain mode~1 and
occur when the frame label returns to zero. A switch, a
velocity reversal and a clock reset are distinguishable events.
\cite{study03}, switch-event tables and log\_summary.json.
\claim{C60}

\section{Energy accounting and renewed clock drive}
The function \filepath{compute_energy} ignores its positions argument.
It returns
\begin{equation}\label{eq:K03}
K^+=\frac12\sum_{i=0}^7\|v_i^+\|_2^2+0.
\end{equation}
The zero is a literal potential-energy placeholder. No mass parameter,
position potential, conservation rule or energy feedback is implemented.
The same global number is written eight times per callback; summing
the column over a block would count it eight times.
\src{S03}{57--61,129,147}. \claim{C61}

At the first $199\to 0$ label reset, $q=200\to 201$, A's K changes
from 0.000922772 to 0.479749699, and B's from 0.000721817 to
0.473368802. The source changes the prescribed gain from
$H(19.9)\simeq0.00261412235$ to $H(0)=4.5$ while retaining state.
Another zero follows, supplying another strong-drive update.
Figure~\ref{fig:F10} shows this chronology for B. These are recorded
drive/damping effects, not spontaneous creation of a physically
conserved energy. \claim{C62}

For either centroid update, let $a_i$ be the source contribution,
$u_i=v_i+a_i$, and $v_i^+=\R(u_i)$. Directly expand the squared norm
and then subtract the damping loss:
\begin{align}
\tfrac12\|u_i\|^2-\tfrac12\|v_i\|^2
 &=v_i\cdot a_i+\tfrac12\|a_i\|^2,\\
L_i&=\tfrac12\bigl(\|u_i\|^2-\|\R(u_i)\|^2\bigr)\geq0,\\
K^+-K&=\sum_i v_i\cdot a_i+\frac12\sum_i\|a_i\|^2-\sum_iL_i.
\label{eq:energybalance}
\end{align}
The inequality follows componentwise from the denominator in RPCO,
which is at least one. This is the accepted source-derived
bookkeeping identity, not an introduced potential or a new dynamical
law. The accepted 31-transition sample reports a balance residual
no greater than $6.67\times10^{-16}$, alongside the independent
velocity-equation checks. \cite{study03}, TIME\_SWITCH\_AND\_MOVEMENT
and equation-residual tables. \claim{C63}

\plain{The diagnostic can rise because the added interaction changes
velocity before damping reduces it. Restarting the H clock restores
a strong input. The source contains no statement that this kinetic
number must remain constant.}

\section{Native movement, coordinate seams and the detector defect}
The position equation separates native increment
$\delta x_i=0.1v_i^+$ from stored coordinate difference
$x_i^+-x_i$. Equation~\eqref{eq:seam} explains their integer-period
offset. In the accepted logs, all checkable native increments are
below 0.062. In contrast, the largest stored jumps are almost ten.
Table~\ref{tab:jumps03} preserves both measurements and their anchors.
\claim{C64}

\begin{table}[htbp]\centering\small
\caption{Original recorded movements, not image-derived estimates.
Anchors use physical CSV lines. Accepted source: S03 log summary
and wrap-event tables.}\label{tab:jumps03}
\begin{tabularx}{\linewidth}{@{}l Y Y@{}}\toprule
Measure & A & B\\\midrule
Maximum native increment &0.061531942; q205, body 2, Step 3, line 1644
&0.060380826; q205, body 5, Step 3, line 1647\\
Maximum stored jump &9.997419049; q1038, body 4, Step 32, line 8310
&9.993505332; q430, body 5, Step 27, line 3447\\
Native increment at that jump &0.012148730&0.015841001\\
Actual seam events &22 body / 22 coordinate events&7 body / 7 coordinate events\\\bottomrule
\end{tabularx}
\end{table}

The source's \filepath{detect_teleport} instead computes
\begin{equation}\label{eq:baddetector}
T_i=\|\W_{10}(x_i^+-x_i)\|_2
\end{equation}
and prints an event if $T_i>1$. This unsigned modulo is neither a
minimal periodic distance nor the native increment. A tiny negative
coordinate difference becomes almost ten. The diagnostic does not
change state, and its flags are not exported in the CSV.
\src{S03}{64--71,132--134}. \claim{C65}

For example, A at q1, body 0, line 10 has coordinate difference
\[
\Delta x=(0.040328749,-0.024777022,0.022525770)
\]
and native length 0.052418691, with no seam. The source diagnostic gives
9.975329933 and would flag it. The accepted analysis reports
8,151 would-flag body events in A and 3,129 in B, including
8,130 and 3,122 events without a seam. It misses one of A's
22 actual seams. These are recomputed diagnostics from original
records; the original console transcript is missing.
\claim{C66}
\EvidenceFig{F11}
\EvidenceFig{F12}

The first B seam is q2, body 2, Step 1, line 20. Z changes from
0.004472603 to 9.965972877 while the native length is 0.054654713.
The corresponding stored jump is 9.961575809.
A's earliest seam occurs at q265, body 3, Step 63, line 2125:
Z changes from 0.001799251 to 9.998477914 with native length
0.005387221. Neither event is a measured native movement of
approximately ten units. \claim{C67}

\section{Body 7 crosses back and forth without a mode change}
In B, body~7 retains mode~1 throughout q400--420. Its X coordinate
crosses the seam at q409,414,418, corresponding to lines
3281,3321,3353. The native lengths are approximately
0.055897106, 0.044266228 and 0.037032498.
Figure~\ref{fig:F13} separates the coordinate jump, VX and mode.
The next mode toggle is q431, outside this displayed window.
\claim{C68}
\EvidenceFig{F13}

At q409 the center crosses from small positive X to a value near ten.
The subsequent wrapped centroid X displacement changes from about
$+4.92$ to $-3.80$. VX then changes from negative to positive,
and the center crosses back at q414. The next displacement becomes
about $+4.93$ and VX turns negative again before q418.
These are accepted measurements of the mixed-state centroid and
velocities, retained in analysis/episode\_records.json.
The arithmetic-mean dependence in Equation~\eqref{eq:centroid}
explains how a coordinate crossing can change the subsequent target.
No wall, collision or added teleport mechanism is present.
\claim{C69}

The observation supports real sign changes in a velocity component
and small native movements through a coordinate seam. It is not
a proof of a periodic orbit, nor an isolated causal experiment
holding every other body fixed. Other centers continue moving and
contribute to the centroid. \claim{C70}

\section{Pair relations and collective interpretations}
The accepted scan of all 28 unordered pairs at stored callbacks
found A's closest pair at q602, Step 199: bodies~2 and~7,
lines 4820/4825, distance 0.111912869. Their velocity cosine is
0.959590042 and relative speed 0.006143850. That is a near,
similarly directed slow pair at a specified instant.
B's closest pair occurs at q118, Step 117: bodies~1 and~4,
lines 947/950, distance 0.813987821, velocity cosine
$-0.568814101$. Its nearest pair is not aligned at that instant.
\cite{study03}, pair-minima tables. \claim{C71}

A narrow descriptive screen in the accepted study used separation
below 0.1, velocity cosine above 0.9 and nonzero speeds above
$10^{-6}$. Neither log met its distance threshold. This is not
a physical threshold, a path-inheritance test or a rejection of
approximate collective following. No route-delay criterion was
imposed on these records. \claim{C72}

\section{The original movie and the candidate correspondence}
The original \filepath{tetrahedral_attractors_logic.mp4} has
200 decoded frames at 20 fps, 1000 by 800 pixels and a duration of
10 seconds. The title, filename, eight tetrahedra, axis range and
viridis/Reds palettes match S03's declarations. Palette family
indicates mode; within either family its phase changes with
$(i+f)\bmod8$. Pale red faces can be difficult to see.
No numeric energy or detector-event overlay is drawn.
\src{S03}{81--97,137--141,158--161}. \claim{C73}
\EvidenceFig{F09}

The accepted inspection decoded every frame and visually reviewed
selected consecutive sheets, an overview, four full-resolution
stills and a comparison panel. A bounded attribution test rendered
\emph{already-recorded} states, never a calculated trajectory.
It compared movie frame 0 with all 18 recorded label-zero states and then
tested frames 0,1,6,34,82,120,184,199 against B callback $q=f+1$.
The fixed geometry, view and palette were not optimized.
\claim{C74}

At movie frame 0, Bq1 gives colored-pixel mask IoU 0.890282 and mean
absolute RGB error 5.516925 on the 0--255 scale. Bq0 gives
0.859051 and 13.001680. Across the eight fixed candidate pairs,
IoU ranges from 0.856495 to 0.890282 and RGB error from
5.516925 to 7.086700. These descriptive image metrics support
a compelling content correspondence, not a calibrated likelihood.
Lossy encoding, transparency and rendering versions affect the
scores. Figure~\ref{fig:F16} selects two of the already-tested
comparisons. \claim{C75}
\EvidenceFig{F16}

\textbf{The formal original-run join remains UNLINKED.}
No export receipt binds a source hash, RNG state, log and movie.
The tested correspondence does not certify all 200 frames or a
unique historical execution. Conditionally, the opening movie
relocation matches Bq1--2's small Z seam crossing, and frame 6
matches Bq7's first mode toggle. No movie observation silently
inherits a CSV measurement as established ground truth.
\claim{C76}

\section{Historical charts and mixed horizontal axes}
Two historical body-7 PNGs are strongly attributed to postprocessing
of B. The position chart's trace shapes and event markers match
the shorter log, but no generating chart script was recovered.
S03 itself contains no PNG export. Figure~\ref{fig:F14} preserves
the historical chart; Figure~\ref{fig:F15} reuses the accepted
diagnostic demonstrating its mixed horizontal coordinates.
\claim{C77}
\EvidenceFig{F14}
\EvidenceFig{F15}

The traces use repeated Step labels 0--199 while markers at
409,414,418 correspond to callback indices of body~7's actual
seam crossings. Connecting physical rows on a repeated Step
axis folds passes together and joins label 199 back to label 0.
The event markers then lie beyond the trace range. This explains
a misleading layout without altering historical bytes or
claiming exact recovery of the lost plotting script.
\claim{C78}

The old GRF/energy chart (Figure~\ref{fig:F17}) has the same
repeated-label issue. Its Energy comes from the global K repeated
on body 7's rows; the title does not make it body 7's individual
kinetic quantity. The three highlighted seam markers also
cannot be equated to the source detector's thousands of
would-flag events. Their selection appears to be different
postprocessing, whose precise algorithm remains missing.
\claim{C79}
\EvidenceFig{F17}
\FloatBarrier
